\documentclass[12pt, a4paper]{article} % Classe: article, report o book [8]

% --- Lingua e Codifica (Supporto Italiano) ---
\usepackage[utf8]{inputenc}  % Codifica input
\usepackage[T1]{fontenc}      % Codifica font
\usepackage[italian]{babel}   % Lingua italiana [2]

% --- Impaginazione e Margini ---
\usepackage[margin=2.5cm]{geometry} % Imposta margini (2.5cm su tutti i lati)

% --- Pacchetti Scientifici/Matematici ---
\usepackage{amsmath, amssymb, amsfonts, amsthm} % Matematica avanzata [2]
\usepackage{graphicx} % Per inserire immagini
\usepackage{hyperref} % Per collegamenti ipertestuali
\usepackage{booktabs} % Per tabelle professionali

% --- Formattazione Testo e Altro ---
\usepackage[document]{ragged2e}
\usepackage{xcolor}    % Per usare i colori 

\title{Frida Lega}
\date{27 Gennaio 2026}


\begin{document}

\maketitle

\section{Ruolo}
Frida Lega è una commerciale con esperienza di più di 10 anni all'interno
dell'azienda e si occupa di :
\begin{itemize}
      \item Supporto ai commerciali junior.
      \item Creazione delle offerte Access per conto di Agenti o rivenditori esterni.
      \item Creazione offerte per clienti diretti.
      \item Formazione interna sull'uso di Access.
      \item Organizzazione e gestione fiere di settore.
\end{itemize}

\section{Situazione attuale}
\subsection{Problemi attuali del sistema}
Dall'intervista a Frida Lega sono emerse le seguenti criticità del sistema
attuale:
\begin{itemize}
      \item \textbf{Lentezza del sistema:} L'attuale database Access risulta lento e con frequenti blocchi.
      \item \textbf{Interfaccia utente datata:} L'aspetto grafico e l'usabilità non sono al passo con le aspettative moderne.
      \item \textbf{Disconnessioni VPN:} Perdita di dati o impossibilità di lavorare efficacemente da remoto a causa di instabilità della connessione.
      \item \textbf{Storico in più database:} Difficoltà nell'accesso e gestione dello storico offerte, attualmente distribuito su più file, uno per anno.
\end{itemize}

\subsection{Funzioni da mantenere}
Nonostante le criticità, Frida ha sottolineato l'importanza di mantenere alcune
funzionalità chiave dell'attuale sistema:
\begin{itemize}
      \item \textbf{Proprietà legate al cliente:} Mantenere la possibilità di associare specifiche tecniche \textit{(es. voltaggio, normative)}
            in base al paese di destinazione del cliente.
      \item \textbf{Gestione clienti e destinazioni:} La struttura a tre livelli (Cliente, Utilizzatore, Destinazione) è considerata efficace.
      \item \textbf{Creazione capitoli}: La struttura a capitoli per le offerte è considerata efficace.
      \item \textbf{Gestione delle scontistiche:} La possibilità di applicare sconti a diversi livelli (singolo componente, capitolo, totale) è fondamentale.
      \item \textbf{Calcoli integrati:} La presenza di calcolatori per potenza o aria pneumatica, anche se da migliorare, è vista come un elemento chiave da mantenere.
            Attualmente viene sfruttata in parte, ma con necessità di interventi manuali di correzione.
      \item \textbf{Traduzioni multilingua:} La capacità di generare offerte in diverse lingue è essenziale per il mercato internazionale.
            Tenere aperta la possibilità di aggiungere ulteriori lingue in futuro \textit{(es. Arabo, Russo, \dots)}.
      \item \textbf{Aggiunta di opzioni alle offerte:} Attualmente c'è la possibilità di gestire dei componenti opzionali che valuterà il cliente
            se accettare senza obbligare l'inserimento forzato nel preventivo.
      \item \textbf{Opzioni sulle stampe:} Possibilità attuale di mostrare o meno i prezzi per prodotto o solo per capitolo o finali.
      \item \textbf{Riferimenti commerciali:} Mantenere la possibilità di compilare offerte per altri commerciali mantenendo traccia del referente originale.
      \item \textbf{Agenti esterni:} Mantenere la possibilità di creare offerte per agenti o rivenditori esterni, loro possono avere solo la vista sulle proprie offerte.
      \item \textbf{Scelta della banca di riferimento:} Attualmente è possibile scegliere la banca di appoggio di riferimento per il cliente. Valutare se può essere guidata nella scelta.
      \item \textbf{Tasto aggiorna prezzi:} Mantenere la funzionalità di aggiornamento massivo dei prezzi in base ai listini caricati,
            è utile mantenere i prezzi congelati al momento della creazione dell'ordine ed aggiornarli solo se richiesto.
      \item \textbf{Campo tipologia offerta:} Mantenere il campo che distingue tra \textit{Ampliamento} e \textit{Ricambi} per facilitare la ricerca futura delle offerte.
\end{itemize}

\subsection{Attuale flusso di lavoro:}
L'attuale processo di creazione delle offerte prevede i seguenti passaggi:
\begin{enumerate}
      \item \textbf{Raccolta requisiti:} Frida si interfaccia con il cliente per comprendere le
            esigenze tecniche e commerciali.
      \item \textbf{Ricerca e copiatura di una vecchia offerta:} Spesso si parte da un'offerta
            precedente simile per velocizzare il processo.
      \item \textbf{Selezione prodotti:} Aggiunta dei prodotti manualmente uno ad uno tramite descrizione.
      \item \textbf{Configurazione offerta:} Creazione dell'offerta suddivisa in capitoli,
            applicando sconti e calcoli necessari.
      \item \textbf{Invio al cliente:} Generazione del documento finale in formato PDF o
            cartaceo ed invio al cliente.
      \item \textbf{Approvazione offerta:} In caso di accettazione, si procede con la
            conferma d'ordine e la comunicazione all'ufficio produzione e l'ufficio amministrativo.
\end{enumerate}

\section{Requisiti funzionali}

\subsection{Gestione offerta e prodotti}
Frida ha evidenziato la necessità di un sistema che faciliti la creazione e
gestione delle offerte commerciali, con particolare attenzione a:
\begin{itemize}
      \item \textbf{Configurazione guidata (Wizard):}
            \begin{itemize}
                  \item Selezione per macro-argomento/macchina \textit{(es. Silos STX)}.
                  \item Apertura automatica di finestre per accessori compatibili (check-box).
                  \item Proposta automatica di ricambi associati alla macro-categoria.
                  \item Alert non bloccanti per componenti mancanti \textit{(es. "Hai messo la
                              tramoggia farina, manca il vibratore")}.
            \end{itemize}
      \item \textbf{Gestione prezzi e sconti:}
            \begin{itemize}
                  \item Scontistica a tre livelli: Singolo componente, Capitolo, Totale offerta.
                  \item Possibilità di visualizzare prezzi netti e di listino.
                  \item Gestione rincari e costi accessori \textit{(es. trasporti, imballi)}.
            \end{itemize}

      \item \textbf{Gestione delle opzioni:}Attualmente è complicata la visualizzazione delle opzioni di vendita di una commessa, per tanto serve:
            \begin{itemize}
                  \item Una migliore gestione visuale delle opzioni.
                  \item Possibilità di raggruppare le opzioni per capitoli.
                  \item Gestione delle opzioni simile a quella della sezione offerta.
            \end{itemize}
      \item \textbf{Gestione sconti:}Possibilità di mantenere la possibilità attuale di applicare rincari o sconti , di questi ultimi si richiede:
            \begin{itemize}
                  \item Sconti a livello di singolo componente.
                  \item Sconti a livello di capitolo.
                  \item Sconti sul prezzo totale offerta senza manodopera.
                  \item Sconto sul prezzo comprensivo di manodopera.
            \end{itemize}
      \item \textbf{Costificazione manodopera:} Attualmente per quantificare la manodopera si richiede direttamnete al reparto addetto una stima in ore,
            da questa si calcola manualmente il costo totale. Frida sottolinea come sia difficile prevedere un sistema automatico in quanto uno stesso silo
            montato in luoghi diversi può richiedere tempi di montaggio molto differenti.
\end{itemize}

\subsection{Dati tecnici e integrazioni}
Il sistema deve integrare dati tecnici e documentazione in modo automatico per
ridurre errori e tempi di inserimento:
\begin{itemize}
      \item \textbf{Codici prodotto:} Utilizzo dei codici del gestionale (Ad Hoc) per allineamento con magazzino e ufficio tecnico.
      \item \textbf{Documentazione allegata:} Link automatico a disegni (Fusion), manuali e certificazioni \textit{(es. conformità, ATEX)} basato sul codice prodotto.
      \item \textbf{Calcolatori integrati:}
            Attualmente esiste un sistema di calcolo , ma non essendo i prodotti aggiornati il conto non è accurato.
            Per migliorare questo aspetto, il nuovo sistema dovrebbe includere:
            \begin{itemize}
                  \item Calcolo automatico preliminare di: Pesi, Volumi, Potenza elettrica, Consumo
                        aria compressa.
                  \item Aggiornamento dinamico dei valori in base ai componenti selezionati.
            \end{itemize}
      \item \textbf{Integrazione calcoli Excel:} Possibilità di integrare calcoli complessi attualmente gestiti tramite file Excel esterni \textit{(es. scale, silos)}
            direttamente all'interno, anche solo come link al file in questione.
      \item \textbf{Stati articoli (EX):} Gestione stati prodotti: \textit{Obsoleto} (non usare), \textit{Sconsigliato} (non standard), \textit{Attivo}.
      \item \textbf{Gestione incoterms:} Scelta guidata per l'inserimento degli incoterms, Frida ha confermato che l'attuale gestione manuale può portare a errori per i meno esperti.
\end{itemize}

\subsection{Gestione Clienti e Anagrafiche}
Il sistema deve gestire una struttura gerarchica a tre livelli, separati ma
collegabili:
\begin{enumerate}
      \item \textbf{Cliente:} Chi acquista/fattura \textit{(es. Italmarco)}.
      \item \textbf{Utilizzatore:} Chi usa l'impianto \textit{(es. Stabilimento produttivo)}.
      \item \textbf{Destinazione:} Luogo fisico di consegna (spesso coincide con l'utilizzatore).
\end{enumerate}
\textit{Requisito specifico:} Le proprietà tecniche \textit{(es. voltaggio, frequenza, normative UL/CSA/ATEX)} devono essere legate al paese di destinazione,
ma sovrascrivibili a livello di singolo cliente.

\section{Workflow e ciclo di vita dell'offerta}

\subsection{Revisioni e storico}
Gestione avanzata delle modifiche post-offerta:
\begin{itemize}
      \item Mantenimento dello stesso numero di offerta con gestione delle
            \textbf{Revisioni} \textit{(es. Offerta 100 Rev. 01, Rev. 02)}.
      \item Storico completo delle modifiche (versioning) per poter recuperare versioni
            precedenti.
      \item Accessibilità ai database storici (almeno ultimi 5-7 anni) per consultazione e
            clonazione offerte.
\end{itemize}

\subsection{Contratto tecnico}
Creazione di un ``doppio binario'' post-vendita:
\begin{itemize}
      \item \textbf{Contratto commerciale:} Quello firmato dal cliente (statico).
      \item \textbf{Contratto tecnico:} Modificabile dai Project Manager per riflettere la realtà costruttiva.
      \item \textbf{Confronto:} Visualizzazione del ``Delta'' (differenze di prezzo/componenti) tra commerciale e tecnico per valutare se emettere Revisione Ordine (RA)
            o assorbire i costi.
\end{itemize}

\section{Requisiti tecnici e infrastrutturali}

\subsection{Offline e sincronizzazione}
Data la necessità di lavorare presso i clienti o in mobilità con connessioni
instabili:
\begin{itemize}
      \item Adozione di un database locale (replica) che si sincronizza con il server (SQL
            Server).
      \item Meccanismo di switch automatico o manuale tra modalità Online/Offline.
      \item Esistenza di un backup locale per evitare perdita dati in caso di caduta VPN.
\end{itemize}

\subsection{Integrazioni esterne}
Il sistema deve avere la possibilità di integrarsi con le seguenti piattaforme
esistenti:
\begin{itemize}
      \item \textbf{Ad Hoc (ERP):} Fonte primaria per anagrafiche articoli e codici.
      \item \textbf{HubSpot (CRM):} Integrazione desiderata per anagrafiche clienti ed evitare doppio inserimento (attualmente l'adozione interna di HubSpot è parziale).
      \item \textbf{File System:} Collegamento a cartelle di rete per recupero disegni tecnici.
\end{itemize}

\section{Note di implementazione}
Frida ha fornito alcune indicazioni utili per la fase di implementazione:
\begin{itemize}
      \item \textbf{Campi obbligatori:} Rendere obbligatoria la distinzione tra ``Ampliamento'' e ``Ricambi'' (attualmente campo discrezionale spesso vuoto)
            per migliorare la ricercabilità futura.
      \item \textbf{Beta tester:} Si consiglia di iniziare la fase di test con figure tecniche/analitiche (Enrico, Thomas) e commerciali esperti (Frida, Natalia).
      \item \textbf{Transizione:} Prevista una fase transitoria con utilizzo parallelo dei due sistemi;
            necessaria migrazione o accessibilità facile ai dati storici.
\end{itemize}

\end{document}